Training-free caching accelerates a video diffusion model by evaluating a subset of the denoising steps and reusing the last exact output elsewhere. Existing methods decide which steps to evaluate with a threshold on a quantity that has already been observed, and spend their design effort on how a skipped step is reconstructed. We measure what a skipped step actually costs, and the cost turns out to be positional. On CogVideoX-2B at thirty steps, skipping the first, second and sixth step costs 0.57, 0.45 and 0.09 LPIPS, while every step after the tenth costs about 0.04, so the first ten steps carry about three quarters of the sensitivity. The profile is a property of the sampler, and two prompts agree on it at 0.94 correlation with identical top eight steps. We present \method, a training-free schedule with an exact head and gaps that grow geometrically along the tail, whose length follows from the measured profile and whose number of calls equals the budget by construction. Evaluated inside one released pipeline at matched calls, \method{} reaches 0.151 LPIPS where the risk-based cache of RA-CFGCache reaches 0.169 on CogVideoX-2B, 0.121 against 0.132 at 49 frames, and 0.134 against 0.154 on Wan2.1, while a threshold-based cache reaches 0.296, 0.290 and 0.199. Code, schedules and all evaluation traces are released at \url{https://github.com/Yizhou-Yang/MPCache}.
